\documentclass[12pt]{article}
\usepackage[utf8]{inputenc}
\usepackage[T1]{fontenc}
\usepackage{hyperref}
\title{Bhutanese Sign Language Language Resources for K-12 Curriculum Development}
\author{Compiled by Claude Sonnet 5.5 (claude-sonnet-5-5) for Andr\'as Kornai}
\date{}
\begin{document}
\maketitle

\section{Introduction}
Bhutanese Sign Language (BhSL, ISO 639-3: dyl) is the indigenous sign language of Bhutan, centred on the Wangsel Institute for the Deaf in Paro. Documentation and standardisation are recent and ongoing.

\section{Primary Educational Resources}
The Ministry of Education's Wangsel Institute page describes BhSL and its research and documentation programme (research from 2002, formal programme from 2015): \url{https://wangsel.education.gov.bt/bhutan-sign-language/}. No published textbooks or primers were found in the searches.

\section{Linguistic and Reference Materials}
\begin{itemize}
\item Wikipedia overview: \url{https://en.wikipedia.org/wiki/Bhutanese_Sign_Language}
\item UNICEF story on deaf students and researchers documenting signs and building a grammar: \url{https://www.unicef.org/bhutan/stories/deaf-students-construct-first-sign-language-bhutan}
\item Deaf cultures and sign languages of the world, Bhutan page: \url{https://theinterpretersfriend.com/indj/dcoew/bhutan.html}
\end{itemize}

\section{Religious and Community Resources}
None identified in the searches.

\section{Digital and Audio Resources}
A BhSL mobile app from the Wangsel Institute (reportedly with a dictionary and quizzes) was reported in Daily Bhutan: \url{https://www.dailybhutan.com/article/breaking-barriers-bhutanese-sign-language-app-for-inclusivity}. A BBS news item also appeared in results: \url{https://bbs.bt/news/?p=195310}. The app itself was not located.

\section{Curriculum Development Context}
Part of government funding for deaf education reportedly goes to developing BhSL as the language of instruction; the first dedicated deaf school (Paro) was approved in 2013. A BhSL-medium K-12 curriculum would need to build on the Wangsel Institute's work.
\end{document}